\SourceLeaf{7}
\begin{center}
THEORY OF THE ATTRACTION OF\par
\bigskip
HOMOGENEOUS ELLIPTIC SPHEROIDAL BODIES\par
\bigskip
TREATED BY A NEW METHOD.\par
\medskip
\rule{.16\textwidth}{.3pt}
\end{center}
\Article{1}
\par
It is well known that the problem of determining exactly the attraction of a homogeneous elliptic spheroidal body on any given point is one of the most difficult problems of physical astronomy, and that several geometers since \textsc{Newton}’s time have applied themselves to it repeatedly and with great energy. At first the investigation was restricted to a spheroid generated by revolving a half ellipse about either axis. The great \textsc{Newton} himself showed how to find the attraction experienced by a point on the axis and at the same time established the relation between the attractions experienced by points inside the spheroid lying on the same diameter (\textit{Princip. Lib.} I. \textit{Prop.} XCI). Next, the keen-minded \textsc{Mac Laurin} used a most elegant synthesis to determine the attraction on points on the spheroid’s surface or in the extended plane of its equator. This also completed the theory of attraction for points inside the spheroid, since \textsc{Newton}’s theorem readily reduced their case to the attraction on points on the surface (De caussa physica fluxus et refluxus maris, in \textit{Recueil des pièces qui ont remporté les prix de l’acad. roi. des sc. T.} IV; \textit{Treatise of fluxions B.} I. \textit{Ch.} 14). The distinguished \textsc{Lagrange} later showed, no less elegantly, how to obtain by analysis what \textsc{Mac Laurin} had worked out by synthesis. Such questions had previously seemed beyond the reach of analysis. His work thus opened the way to further advances (\textit{Nouv. Mém. de l’Acad. de Berlin} 1773). What still remained to be worked out was the attraction on points outside the spheroid lying on neither the extension of the axis nor that of the equator,
